\chapter{第3章\texorpdfstring{\hspace{0.5em}}{ }刚体力学基础}

前两章把研究对象简化为质点。实际上物体的机械运动除平动外，尚有转动和振动。本章把研究对象从质点扩大为由无数连续分布的质点组成、有形状和大小的\textbf{刚体}（rigid body）。所谓刚体，是指在任何外力作用下，其内部各质点间的相对位置均保持不变的物体，即形变可以忽略的理想模型。本章先引入描述物体转动的重要物理量——角动量和转动惯量，进而讨论引起角动量变化的三个规律：力矩的瞬时作用规律（刚体定轴转动定律）、力矩对时间的累积作用规律（角动量定理和角动量守恒定律）、力矩对空间的累积作用规律（刚体定轴转动的动能定理和机械能守恒定律），最后介绍刚体的平面平行运动。重点是刚体的定轴转动。

\section{刚体运动的基本形式}
    \subsection{平动和转动}
    刚体可以看成是由无数个连续分布的质点所组成的质点系，每个质点称为刚体的一个\textbf{质量元}。刚体在运动中各部分相对位置均保持不变，因此只要知道刚体上某一点的运动情况，则刚体上任一点乃至刚体运动的全貌便可确定。

    刚体的运动可以分为平动和转动，刚体任何复杂的运动均可看成是二者的叠加。
    \begin{enumerate}
        \item \textbf{平动}（translation）：刚体上任一给定直线（或任意两质点间连线）在运动中空间方向始终不变而保持平行，这样的运动称为平动，如活塞、抽屉的运动。刚体平动时，刚体上各点的运动轨迹均相同，同一时刻各点具有相同的速度和加速度，因此完全可以简化为质点处理，通常用刚体质心的运动代表整个刚体的平动。
        \item \textbf{转动}（rotation）：如果刚体上所有质点都绕同一直线作圆周运动，这种运动称为刚体的转动，这条直线称为\textbf{转轴}（rotation axis）。
        \begin{itemize}
            \item 若整个转轴相对参考系静止，称为\textbf{定轴转动}（如门窗、车床工件的转动）；
            \item 若转轴上只有一点相对参考系静止而转轴方向不断变动，称为\textbf{定点转动}（如玩具陀螺、厢顶电风扇的转动）。
        \end{itemize}
    \end{enumerate}
    车轮的滚动总可分解为车轮转轴的平动和整个车轮绕转轴的转动。刚体的定轴转动及其规律是本章的重点，也是研究整个刚体力学的基础。

    \subsection{描述刚体转动的物理量}
    刚体定轴转动时具有两个特征：轴上各点始终静止不动；轴外各质点虽有各自与轴垂直的转动平面、到轴的距离（转动半径）不同、相同时间内转过的线位移不同，但转过的\textbf{角位移相同}。因此可在刚体上任意选定一点，通过研究该点绕定轴的转动来描述整个刚体的定轴转动。

    通常选取与轴相固联的惯性系为参考系，以相对参考系的任一固定方向 $Ox$ 作为参考方向，质点对转轴交点 $O$ 的位矢 $\boldsymbol{r}$ 与该方向的夹角 $\theta$ 即为该质点的\textbf{角坐标}。刚体转动时
    \keyformula{\omega=\frac{\mathrm{d}\theta}{\mathrm{d}t},\qquad
    \alpha=\frac{\mathrm{d}\omega}{\mathrm{d}t}=\frac{\mathrm{d}^{2}\theta}{\mathrm{d}t^{2}}}
    角速度的方向沿转轴，遵从右手螺旋定则；$\omega$ 变大时 $\alpha$ 与 $\omega$ 同向，$\omega$ 变小时二者反向。定轴转动中角速度、角加速度只有沿转轴的两个方向，计算中常作标量处理。

    质点线速度与加速度同角量的关系为
    \keyformula{\boldsymbol{v}=\boldsymbol{\omega}\times\boldsymbol{r},\qquad
    \boldsymbol{a}=\boldsymbol{\alpha}\times\boldsymbol{r}+\boldsymbol{\omega}\times\boldsymbol{v}}
    其中 $\boldsymbol{\alpha}\times\boldsymbol{r}$ 是切向加速度，大小 $a_{\mathrm{t}}=r\alpha$；$\boldsymbol{\omega}\times\boldsymbol{v}$ 是法向加速度，大小 $a_{\mathrm{n}}=r\omega^{2}=v^{2}/r$。

    \begin{center}
    \begin{tabular}{lll}
    \toprule
    运动学量 & 平动 & 定轴转动 \\
    \midrule
    位置 & $x$ & $\theta$ \\
    速度 & $v=\mathrm{d}x/\mathrm{d}t$ & $\omega=\mathrm{d}\theta/\mathrm{d}t$ \\
    加速度 & $a=\mathrm{d}v/\mathrm{d}t$ & $\alpha=\mathrm{d}\omega/\mathrm{d}t$ \\
    匀变速公式 & $v=v_{0}+at$ & $\omega=\omega_{0}+\alpha t$ \\
    匀变速公式 & $x=x_{0}+v_{0}t+\frac{1}{2}at^{2}$ & $\theta=\theta_{0}+\omega_{0}t+\frac{1}{2}\alpha t^{2}$ \\
    匀变速公式 & $v^{2}-v_{0}^{2}=2a(x-x_{0})$ & $\omega^{2}-\omega_{0}^{2}=2\alpha(\theta-\theta_{0})$ \\
    \bottomrule
    \end{tabular}
    \end{center}

\section{定轴转动刚体的角动量\texorpdfstring{\hspace{0.5em}}{ }转动惯量}
    \subsection{刚体对定轴的角动量}
    当刚体绕定轴转动时，刚体上各质元某一瞬时均以相同的角速度 $\boldsymbol{\omega}$ 绕该轴作圆周运动。设刚体上某质元 $\Delta m_{i}$ 距轴的距离为 $r_{i}$，其速度方向与位矢垂直，故对该轴的角动量大小为 $L_{i}=\Delta m_{i}v_{i}r_{i}=\Delta m_{i}r_{i}^{2}\omega$。整个刚体对该轴的角动量等于各质元角动量的总和：
    \keyformula{L_{z}=\sum_{i}L_{i}=\left(\sum_{i}\Delta m_{i}r_{i}^{2}\right)\omega}
    括号中的量只与刚体的质量分布和转轴位置有关，记为转动惯量 $J$，于是
    \keyformula{L_{z}=J\omega}
    刚体对转轴的角动量方向与角速度方向一致，计算中通常先设定正方向，再用正负号表示 $L_{z}$ 的方向。

    \subsection{转动惯量}
        \subsubsection{定义与物理意义}
        \textbf{转动惯量}（moment of inertia）：质点系内每个质点的质量与该质点到转轴距离平方之积的总和。
        \keyformula{J=\sum_{i}\Delta m_{i}r_{i}^{2}\ \longrightarrow\ J=\int_{m}r^{2}\,\mathrm{d}m}
        对质量连续分布的刚体用积分（$\mathrm{d}m=\lambda\,\mathrm{d}l$、$\sigma\,\mathrm{d}S$ 或 $\rho\,\mathrm{d}V$）；对质量面分布或体分布的物体，$r$ 恒为质元到转轴的垂直距离。SI 单位 $\mathrm{kg\cdot m^{2}}$，量纲 $\mathrm{ML^{2}}$。

        转动惯量的大小取决于：
        \begin{enumerate}
            \item 刚体的\textbf{总质量}；
            \item 质量相对于给定轴的\textbf{分布}——总质量一定时，质量分布离轴越远，转动惯量越大；
            \item \textbf{转轴的位置}——同一刚体，转轴不同，转动惯量不同。
        \end{enumerate}
        因此说“某刚体的转动惯量”时必须指明是对哪一条轴。转动惯量在转动中的地位相当于质量在平动中的地位，它量度刚体转动惯性的大小。

        \subsubsection{回转半径}
        设物体总质量 $m=\sum\Delta m_{i}$，对给定轴的转动惯量为 $J$，定义对该转轴的\textbf{回转半径}（radius of gyration）$r_{G}$ 为
        \keyformula{r_{G}=\sqrt{\frac{J}{m}},\qquad J=mr_{G}^{2}}
        即把全部质量集中在距轴 $r_{G}$ 处时，其转动惯量与原刚体相同。

        \subsubsection{平行轴定理}
        若质量为 $m$ 的刚体对过质心 $C$ 的轴的转动惯量为 $J_{C}$，则对与该轴平行、相距为 $d$ 的另一轴的转动惯量为
        \keyformula{J=J_{C}+md^{2}}
        可见在一切平行轴中，过质心的轴对应的转动惯量最小。

        \subsubsection{正交轴定理}
        薄板形刚体对板面内两条相互正交轴的转动惯量之和，等于它对过两轴交点且垂直于板面的那条轴的转动惯量：
        \keyformula{J_{z}=J_{x}+J_{y}}
        该定理只适用于薄板（面分布）刚体。

        \subsubsection{常见刚体的转动惯量}
        \begin{center}
        \begin{tabular}{ll}
        \toprule
        刚体（质量为 $m$） & 转轴与转动惯量 \\
        \midrule
        细杆（长 $l$） & 过中心且垂直于杆：$J=\dfrac{1}{12}ml^{2}$ \\
        细杆（长 $l$） & 过一端且垂直于杆：$J=\dfrac{1}{3}ml^{2}$ \\
        圆环、薄圆筒（半径 $R$） & 过中心轴：$J=mR^{2}$ \\
        圆盘、实心圆柱（半径 $R$） & 过中心轴：$J=\dfrac{1}{2}mR^{2}$ \\
        球壳（半径 $R$） & 过直径：$J=\dfrac{2}{3}mR^{2}$ \\
        实心球（半径 $R$） & 过直径：$J=\dfrac{2}{5}mR^{2}$ \\
        矩形薄板（边长 $a,b$） & 过中心且垂直板面：$J=\dfrac{1}{12}m(a^{2}+b^{2})$ \\
        \bottomrule
        \end{tabular}
        \end{center}
        \corepoint{记忆要点：质量越向外分布，系数越大（圆环 $1$ $>$ 圆盘 $1/2$ $>$ 实心球 $2/5$）；细杆过端点是过中心的 $4$ 倍，可由平行轴定理 $J=ml^{2}/12+m(l/2)^{2}$ 验证。}

\section{刚体定轴转动定律}
    \subsection{刚体角动量的时间变化率\texorpdfstring{\hspace{0.5em}}{ }对转轴的力矩}
    质点系角动量对时间的变化率等于系统所受的合外力矩，即 $\dfrac{\mathrm{d}\boldsymbol{L}}{\mathrm{d}t}=\boldsymbol{M}$。

    力作用于定轴转动的刚体上时，外力 $\boldsymbol{F}$ 可分解为平行于转轴的分量 $\boldsymbol{F}_{\parallel}$ 和垂直于转轴（在转动平面内）的分量 $\boldsymbol{F}_{\perp}$。$\boldsymbol{F}_{\parallel}$ 的力矩垂直于转轴，对绕该轴的转动不起作用；只有转动平面内的 $\boldsymbol{F}_{\perp}$ 的力矩才起作用。于是定轴转动刚体角动量的时间变化率与外力对转轴的力矩的关系为
    \keyformula{\frac{\mathrm{d}\boldsymbol{L}}{\mathrm{d}t}=\boldsymbol{M}_{z}=\boldsymbol{r}\times\boldsymbol{F}_{\perp}}
    力矩的大小 $M_{z}=rF_{\perp}\sin\varphi=F_{\perp}d$，$d$ 为力臂。力矩表征力对物体转动的作用，是力的大小、方向及作用点诸因素共同决定的物理量；单位 $\mathrm{N\cdot m}$。

    \subsection{刚体定轴转动定律}
    \textbf{刚体定轴转动定律}：刚体所受的对于某定轴的合外力矩，等于刚体对该定轴的转动惯量与刚体在此合外力矩作用下所获得的角加速度的乘积。
    \keyformula{M_{z}=J\alpha}
    该式是合外力矩对刚体的瞬时作用规律，式中各量必须是对同一时刻、同一刚体、同一转轴而言。
    \corepoint{定轴转动定律与牛顿第二定律地位相当：$M_{z}\leftrightarrow F$，$J\leftrightarrow m$，$\alpha\leftrightarrow a$。解题时定轴转动定律与牛顿第二定律（对参与运动的平动物体）联立，再由绳与轮之间“不打滑”的运动学关系 $a=R\alpha$ 补充一个方程。}

\section{刚体的角动量定理和角动量守恒定律}
    \subsection{刚体的角动量定理}
    刚体可看作由质元组成的质点系，其角动量定理的微分形式与积分形式为
    \keyformula{\boldsymbol{M}=\frac{\mathrm{d}\boldsymbol{L}}{\mathrm{d}t},\qquad
    \int_{t_{1}}^{t_{2}}\boldsymbol{M}\,\mathrm{d}t=\boldsymbol{L}_{2}-\boldsymbol{L}_{1}}
    对定轴转动的刚体，各矢量均可沿轴用标量表示：
    \keyformula{\int_{t_{1}}^{t_{2}}M_{z}\,\mathrm{d}t=L_{z2}-L_{z1}=J\omega_{2}-J\omega_{1}}
    即在一段时间内作用在刚体上的外力的\textbf{冲量矩}等于刚体在同一段时间内角动量的增量。

    \subsection{刚体的角动量守恒定律}
    若刚体（或系统）所受对某定轴的合外力矩为零，则对该轴的角动量保持不变：
    \keyformula{J\boldsymbol{\omega}=\text{常矢量}}
    三种常见情形：
    \begin{enumerate}
        \item 对\textbf{定轴转动的刚体}，转动过程中 $J$ 不变，合外力矩为零时刚体以恒定角速度绕轴转动。轮船、飞机、火箭上用作导航定向的回转仪即利用此原理。
        \item 对\textbf{定轴转动的非刚体}，各质元到轴的距离可变，$J$ 也可变。合外力矩为零时 $J\omega=$ 常量，$J$ 与 $\omega$ 成反比：人站在可绕竖直轴转动的转台上，两臂平伸、手握重物，转动后收拢双臂转速加快，张开双臂转速变慢；滑冰运动员、跳水运动员、芭蕾演员常借此完成动作。
        \item 对\textbf{物体系统}，当若干质点和刚体组成系统且对同一轴的外力矩之和为零时，系统对该轴的总角动量守恒：
        \[
        \sum_{i}J_{i}\boldsymbol{\omega}_{i}+\sum_{j}\left(\boldsymbol{r}_{j}\times m_{j}\boldsymbol{v}_{j}\right)
        =\text{常矢量}
        \]
    \end{enumerate}
    角动量守恒定律与动量守恒定律、能量守恒定律一样，是自然界的普遍规律；在微观领域也常用它推算粒子衰变产物的运动状态。

    \subsection{进动}
    当陀螺不旋转时，受重力矩作用会立即倾倒；但当它绕自身对称轴高速自旋时，尽管同样受重力矩作用却不会倒下。陀螺高速自转的同时，其自身对称轴还绕竖直轴作回旋运动，这种现象称为\textbf{进动}（precession），又称回转运动。

    设陀螺质量为 $m$、对自身轴的转动惯量为 $J$、自旋角速度为 $\omega$，则自旋角动量
    \[
    \boldsymbol{L}=J\boldsymbol{\omega}
    \]
    方向沿自身对称轴。以陀螺与地面的接触点 $O$ 为参考点，重力矩为
    \[
    \boldsymbol{M}=\boldsymbol{r}_{c}\times m\boldsymbol{g}
    \]
    其方向与 $\boldsymbol{L}$ 垂直，故只能改变 $\boldsymbol{L}$ 的方向而不改变其大小。经 $\mathrm{d}t$ 时间，角动量增量为 $\mathrm{d}\boldsymbol{L}=\boldsymbol{M}\,\mathrm{d}t$，方向与 $\boldsymbol{M}$ 平行，于是 $\boldsymbol{L}$ 在水平面内持续偏转，自转轴绕竖直轴扫出一个圆锥面。设自转轴与圆锥轴线夹角为 $\theta$，则
    \[
    \mathrm{d}\varphi=\frac{|\mathrm{d}\boldsymbol{L}|}{L\sin\theta}
    \]
    相应的\textbf{进动角速度}大小为
    \keyformula{\Omega=\frac{\mathrm{d}\varphi}{\mathrm{d}t}
    =\frac{M}{L\sin\theta}=\frac{M}{J\omega\sin\theta}}
    可见进动角速度与外力矩成正比，与自旋角动量（从而与自旋角速度）成反比。

    \corepoint{该进动公式是近似关系：推导中把陀螺绕自身轴的自旋角动量当作总角动量，忽略了进动角动量，因此只有 $\omega\gg\Omega$ 时才成立。一般情况下自转轴还会周期性摆动，称为章动。}

    具有自身对称轴、能产生进动现象的物体系统称为\textbf{回转仪}。进动效应在工程上应用广泛：炮膛内壁刻出螺旋线（来复线），使炮弹出膛后绕自身轴高速自旋，空气阻力矩便只使其轴线绕前进方向进动而不再翻转；地球本身是一个巨大的回转仪，在太阳及其他星体引力作用下也会产生进动；微观上，绕核运动的电子处在外磁场中也会以外磁场方向为轴进动。

\section{刚体定轴转动的动能定理和机械能守恒定律}
    \subsection{刚体转动动能和势能}
    设刚体作定轴转动，某时刻角速度大小为 $\omega$，刚体上任一质元 $\Delta m_{i}$ 距轴 $r_{i}$，其动能为
    \[
    \Delta E_{\mathrm{k}i}=\frac{1}{2}\Delta m_{i}v_{i}^{2}=\frac{1}{2}\Delta m_{i}r_{i}^{2}\omega^{2}
    \]
    整个刚体的转动动能为各质元动能之和：
    \keyformula{E_{\mathrm{k}}=\sum_{i}\Delta E_{\mathrm{k}i}
    =\frac{1}{2}\left(\sum_{i}\Delta m_{i}r_{i}^{2}\right)\omega^{2}
    =\frac{1}{2}J\omega^{2}}
    刚体的重力势能等于各质元重力势能之和。设质心相对势能零点的高度为
    $h_{C}=\dfrac{\sum\Delta m_{i}h_{i}}{m}$，则
    \keyformula{E_{\mathrm{p}}=\sum_{i}\Delta m_{i}gh_{i}=mgh_{C}}
    即刚体的重力势能相当于全部质量集中在质心处的质点所具有的重力势能。

    \subsection{力矩的功}
    设刚体在外力作用下有元角位移 $\mathrm{d}\theta=\omega\,\mathrm{d}t$，力 $\boldsymbol{F}_{i}$ 作用点的元位移 $\mathrm{d}\boldsymbol{l}_{i}=\boldsymbol{v}_{i}\,\mathrm{d}t$，且 $\boldsymbol{v}_{i}=\boldsymbol{\omega}\times\boldsymbol{r}_{i}$，则元功
    \[
    \mathrm{d}W_{i}=\boldsymbol{F}_{i}\cdot\left(\boldsymbol{\omega}\times\boldsymbol{r}_{i}\right)\mathrm{d}t
    \]
    由矢量混合积的轮换性 $\boldsymbol{a}\cdot(\boldsymbol{b}\times\boldsymbol{c})=\boldsymbol{b}\cdot(\boldsymbol{c}\times\boldsymbol{a})$，得
    \[
    \mathrm{d}W_{i}=\boldsymbol{\omega}\cdot\left(\boldsymbol{r}_{i}\times\boldsymbol{F}_{i}\right)\mathrm{d}t
    \]
    所有外力的元功之和为
    \keyformula{\mathrm{d}W=\sum_{i}\boldsymbol{\omega}\cdot\boldsymbol{M}_{i}\,\mathrm{d}t
    =\boldsymbol{M}_{\text{外}}\cdot\mathrm{d}\boldsymbol{\theta}}
    即力对刚体做的功可用力矩与角位移的乘积计算，故称\textbf{力矩的功}。当角位置由 $\theta_{1}$ 变到 $\theta_{2}$ 时
    \[
    W=\int_{\theta_{1}}^{\theta_{2}}M_{\text{外}}\,\mathrm{d}\theta
    \]
    力矩的功率为
    \keyformula{P=\frac{\mathrm{d}W}{\mathrm{d}t}=M_{\text{外}}\omega}

    \subsection{刚体定轴转动的动能定理}
    由力矩的元功表达式及转动定律 $M=J\alpha$
    \[
    \mathrm{d}W=M\,\mathrm{d}\theta=J\alpha\,\mathrm{d}\theta
    =J\frac{\mathrm{d}\omega}{\mathrm{d}t}\mathrm{d}\theta=J\omega\,\mathrm{d}\omega
    \]
    从 $t_{1}$（$\theta_{1},\omega_{1}$）积分到 $t_{2}$（$\theta_{2},\omega_{2}$），得
    \keyformula{W=\int_{\theta_{1}}^{\theta_{2}}M\,\mathrm{d}\theta
    =\frac{1}{2}J\omega_{2}^{2}-\frac{1}{2}J\omega_{1}^{2}}
    即刚体转动动能的增量等于合外力矩对刚体所做的功。

    \subsection{刚体定轴转动的机械能守恒定律}
    若刚体除受外力矩外还受摩擦力矩和重力矩作用，按动能定理有
    \[
    W=\int_{\theta_{1}}^{\theta_{2}}\left(M_{\text{外}}+M_{\text{非保内}}+M_{\text{重}}\right)\mathrm{d}\theta
    =\left(\frac{1}{2}J\omega_{2}^{2}+mgh_{C2}\right)-\left(\frac{1}{2}J\omega_{1}^{2}+mgh_{C1}\right)
    \]
    这就是重力场中刚体定轴转动的功能原理。若刚体只受保守力矩作用，则机械能守恒：
    \keyformula{\frac{1}{2}J\omega^{2}+mgh_{C}=\text{常量}}
    对包含刚体的力学系统，只有在保守内力做功、外力与非保守内力（矩）做功均为零时，系统的机械能（转动动能 + 平动动能 + 重力势能）才守恒。
    \corepoint{刚体的机械能守恒是“转动动能、质心平动动能、重力势能”三者的总和守恒；有摩擦或有绳与轮间的滑动时，必须用功能原理而不能直接用机械能守恒。}

\section{刚体的平面平行运动}
    \subsection{刚体平面平行运动的描述}
    若刚体上任一质元的运动轨迹都平行于某一固定平面，这种运动称为\textbf{平面平行运动}。其特点是：每个质元的轨迹都是平面曲线且各平面互相平行；刚体在运动中转轴始终保持平行且垂直于该固定平面。

    研究方法是把运动分解为\textbf{随质心的平动}与\textbf{绕过质心并垂直于运动平面的转轴的转动}。平动部分用质心速度 $\boldsymbol{v}_{C}$、质心加速度 $\boldsymbol{a}_{C}$ 描述，转动部分用角速度 $\boldsymbol{\omega}$、角加速度 $\boldsymbol{\alpha}$ 描述。由运动叠加原理，刚体上任一点（相对质心的位矢为 $\boldsymbol{r}$）的速度与加速度为
    \keyformula{\boldsymbol{v}=\boldsymbol{v}_{C}+\boldsymbol{\omega}\times\boldsymbol{r},\qquad
    \boldsymbol{a}=\boldsymbol{a}_{C}+\boldsymbol{a}_{\text{绕}C}}

    \subsection{刚体平面平行运动的动力学规律}
    随质心的平动遵从质心运动定理
    \keyformula{\sum_{i}\boldsymbol{F}_{i}=m\boldsymbol{a}_{C}
    =m\frac{\mathrm{d}\boldsymbol{v}_{C}}{\mathrm{d}t}}
    其分量式为 $\sum F_{xi}=ma_{Cx}$、$\sum F_{yi}=ma_{Cy}$，其中 $m$ 为刚体总质量，各外力不一定通过质心。

    绕通过质心的转轴的转动遵从定轴转动定律
    \keyformula{M_{C}=J_{C}\alpha}
    这两个方程就是求解刚体平面平行运动动力学问题的基本方程组。刚体的动能也由两部分组成：
    \keyformula{E_{\mathrm{k}}=\frac{1}{2}mv_{C}^{2}+\frac{1}{2}J_{C}\omega^{2}}

    \subsection{纯滚动}
    若刚体在滚动过程中与接触面之间无相对滑动，这种运动称为\textbf{纯滚动}。圆柱（或圆球）每滚过一周，中心轴前进的距离等于横截面的周长，于是中心轴前进距离 $x_{C}$ 与转角 $\theta$ 的关系为 $x_{C}=R\theta$。逐次求导得\textbf{纯滚动条件}：
    \keyformula{x_{C}=R\theta,\qquad v_{C}=R\omega,\qquad a_{C}=R\alpha}
    三式中求解问题时只需用到其中一个。纯滚动时接触点为\textbf{瞬时转动中心}，接触点的速度为零（因此静摩擦力不做功），故纯滚动的机械能守恒。

    通常联立质心运动定理、绕质心转动定律与纯滚动条件，即可求解纯滚动问题。
    \corepoint{判断“纯滚动”与“又滑又滚”的差别就在于 $v_{C}=R\omega$ 是否成立：前者静摩擦力不做功、机械能守恒；后者滑动摩擦力做负功、机械能不守恒。}
